% TeX_root = ../main.tex

\marginnote{\scriptsize 28/01/2025 }

\begin{example}
  Let $(X, \Sigma, \mu)$ be a measure space. Define $E : \Sigma \to
  B(L^{2}(X, \mu)):= A \to M_{\chi_A}$. It is easy to see that $E$
  satisfies the first 3 properties of a spectral measure.
  \textcolor{red}{verify the 4th}.
\end{example}

\begin{proposition}
  \label{prop:spectral_measure_gives_complex_measure}
  Let $E$ be a spectral measure on $(X, \Sigma, \mathcal{H})$. Then
  for every $\xi, \eta \in \mathcal{H}$.
  \begin{align*}
    E_{\xi, \eta}(A) = \langle  E(A) \xi ,  \eta \rangle
  \end{align*}
  defines a finite measure on $(X, \Sigma)$ with $\|E_{\xi, \eta}\|
  \le \|\xi\| \|\eta\|$.
\end{proposition}
\begin{proof}
  \textcolor{red}{verify}
\end{proof}

\begin{lemma}
  Let $\{ P_i \}_{i \in I}$ be a family of pairwise orthogonal
  projections on a Hilbert space $\mathcal{H}$. Then there exist a
  unique projection $P$ on $\mathcal{H}$ such that $\forall \xi, \eta
  \in \mathcal{H}$,
  \begin{align*}
    \sum_{i \in I} \langle P_i \xi ,  \eta \rangle = \langle P \xi ,
    \eta \rangle
  \end{align*}
\end{lemma}
This means the convergence of orthogonal projections is not in the
norm sense, but rather in the sense above. For example
\begin{example}
  Consider $P_n = P_{\delta_n}$ in $B(\ell^{2}(\mathbb{N}))$, and let
  $  Q_n = \sum_{i = 1}^{n} P_i$. Then clearly $Q_i$ doesn't converge
  in norm, but rather in the above sense to the identity map in
  $B(\ell^{2}(\mathbb{N}))$.
\end{example}

\begin{lemma}
  \label{lem:resiz_representation_for_sesquilinear}
  Let $B(\cdot, \cdot)$ be a bounded sesquilinear form on a Hilbert
  space $\mathcal{H}$. Then there exist a unique $T \in
  B(\mathcal{H})$ such that for all $\xi, \eta \in \mathcal{H}$,
  \begin{align*}
    B(\xi, \eta) = \langle T \xi ,  \eta \rangle
  \end{align*}
\end{lemma}

\begin{proposition}
  \label{prop:spectral_measure_and_integration}
  Let $\phi$ be a bounded linear function on $(X, \Sigma)$. Then
  there exist a unique $  T_\phi \in B(\mathcal{H})$ such that for
  all $\xi, \eta \in \mathcal{H}$
  \begin{align*}
    \int_X \phi \ d  E_{\xi, \eta} = \langle T_\phi \xi ,  \eta \rangle
  \end{align*}
  and we denote $T_\phi:= \int_X \phi \ d  E$
\end{proposition}
\begin{proof}
  See that the integral is sesquilinear on $\xi, \eta$. Then use
  \autoref{lem:resiz_representation_for_sesquilinear} with
  \autoref{prop:spectral_measure_gives_complex_measure}.
\end{proof}

Note that the set $\mathcal{M}(X)$ of all bounded measurable functions on $X$,
equipped with sup norm is a Banach space. Moreover, with pointwise
product and complex conjugation it turns into a commutative C$^*$ algebra.

\begin{theorem}
  Let $(X, \Sigma)$ be a measure space and $\phi$ be a bounded linear
  function. The map
  \begin{align*}
    \mathcal{M}(X) \to B(\mathcal{H}):= \phi \mapsto \int_ X \phi \ d E
  \end{align*}
  is a contractive $*$-homomorphism.
\end{theorem}
\begin{proof}
  To show that it is a $*$-homomorphism, observe that
  \begin{align*}
    \overline{E_{\xi, \eta}} = E_{\eta, \xi}
  \end{align*}
  and then proceed. \textcolor{red}{verify}

  To show that the map is multiplicative, we need to show that
  \begin{align*}
    \int \phi \psi \ d  E = \int \phi \ d  E \circ \int \psi \ d  E
  \end{align*}
  See the picture and use \label{prop:spectral_measure_and_integration}.
\end{proof}

\begin{theorem}
  Let $X$ be a compact Hausdorff space and $\mathcal{H}$ be a
  Hilbert space, and $\Psi: C(X) \to B(\mathcal{H})$ is a
  $*$-homomorphism (representation). Then there exist a unique
  spectral measure on $(X, \mathcal{B}, \mathcal{H})$ ($\mathcal{B}$
  being the Borel sigma algebra on $X$), such that
  \begin{align*}
    \Psi(f) = \int_X  f \ d E
  \end{align*}
\end{theorem}

% \begin{definition}
%   A C$^*$ algebra is a Banach *-algebra satifying the norm equality
%   $\|a^*a\| = \|a\|^2$, for all $a \in \mathcal{A}$.
% \end{definition}
